\section{Task 2}

\subsection{Overview of QR decomposition}

QR decomposition factors a matrix \( A \) into the product of two matrices:
\[
A = QR,
\]  
where \( Q \) is orthogonal (\( Q^T Q = I \)) and \( R \) is upper triangular. The method is widely used in numerical computation, particularly for computing matrix eigenvalues.
Whereas the power method repeatedly multiplies \( A \) by a vector and normalizes the result, the QR algorithm operates on a complete vector basis. QR decomposition leads to basis orthogonalization and vector normalization across successive iterations.

\subsection{Results and conclusions:}

Numerical-error results for QR decomposition are as follows:

\textbf{Error definitions:}
\begin{itemize}
\item \textbf{Error 1:} Checks the unitarity of \( Q \), that is, \( Q^T Q = I \). For an ideal QR decomposition, \( Q \) should be orthogonal.
\item \textbf{Error 2:} Checks whether the product of \( Q \) and \( R \) reconstructs the original matrix \( A \), that is, \( Q R = A \). This verifies the decomposition itself.
\item \textbf{Error 3:} Checks the unitarity of \( Q1 \), obtained by further transformations of \( Q \). For orthogonal matrices, \( Q1 \) should also be orthogonal.
\item \textbf{Error 4:} Checks whether the product of \( Q1 \) and \( R1 \) reconstructs \( A \). As for error 2, this verifies the decomposition after additional transformations.
\end{itemize}
The table gives numerical errors for different matrix sizes:

\begin{center}
\begin{tabular}{|l|l|l|l|l|}
\hline 
$n$   & error-1       & error-2       & error-3       & error-4       \\ \hline
$3$   & $8.8817e-16$ & $5.5512e-17$ & $4.4408e-16$ & $3.3308e-16$ \\ \hline
$500$ & $5.4369e-13$ & $3.5417e-14$ & $1.1103e-15$ & $7.9382e-15$ \\ \hline
$100$ & $3.9967e-15$ & $5.6844e-14$ & $1.7765e-15$ & $3.1265e-13$ \\ \hline
\end{tabular}
\end{center}

\textbf{Conclusions:}
\begin{itemize}
    \item For small matrices (\( n = 3 \)), numerical errors are very small, indicating high precision of the QR method at small matrix sizes.
    \item As matrix size increases (\( n = 500 \)), errors grow. This may reflect accumulated rounding errors and limited numerical precision, which become more visible for large matrices. The errors remain small but are significantly larger.
    \item For the symmetric matrices considered in the last table row, errors decrease, suggesting better numerical stability of the QR method for these matrices.
    \item Overall, the QR method has good numerical stability, particularly for symmetric matrices, where convergence is clearer and errors are smaller.
\end{itemize}
